\chapter{The Filesystem}\label{ch:sh-fs}

Unix presents everything as a file in one tree: your documents, your keyboard,
your disks, the running kernel's own state. That claim is repeated so often it
has stopped meaning anything, so here is the version that matters in practice:
there is exactly one namespace, it starts at \file{/}, and the commands in this
chapter are how you move through it.

\section{Paths}\label{sec:sh-paths}

\begin{definition}[Absolute and relative paths]
	An \vocab{absolute} path begins at the root of the tree and so starts with
	\file{/}. A \vocab{relative} path begins wherever you currently are, and does
	not.
\end{definition}

\begin{keys}[0.18]
	\file{/}          & The root of the whole tree                                     \\
	\file{.}          & The current directory                                          \\
	\file{..}         & The parent directory                                           \\
	\file{\textasciitilde}   & Your home directory; \file{\textasciitilde konyi} is someone else's \\
	\file{-}          & As an argument to \cmd{cd}, the previous directory             \\
\end{keys}

\begin{eg}
	From \file{/home/konyi/Documents}, the file \file{/home/konyi/Documents/LaTeX/CSTools.tex}
	can be named \file{LaTeX/CSTools.tex}, \file{./LaTeX/CSTools.tex}, or
	\file{\textasciitilde/Documents/LaTeX/CSTools.tex}. All three are the same
	file; which you use is a question of what you want to stay true when you move.
\end{eg}

The parts of a path have names, because commands use them:

\begin{shell}
$ basename /home/konyi/notes.tex     # the last component
notes.tex
$ dirname /home/konyi/notes.tex      # everything before it
/home/konyi
$ realpath ../Documents              # resolve to an absolute path
/home/konyi/Documents
\end{shell}

\begin{note}
	A leading dot makes a file \vocab{hidden}: it is not listed by \cmd{ls} and not
	matched by \texttt{*}. This is not a permission or a security feature --- it is
	a naming convention that \cmd{ls} honours, which is why your home directory
	looks empty until you run \cmd{ls -a}.
\end{note}

\section{Navigating}\label{sec:sh-navigate}

\begin{keys}[0.22]
	\cmd{pwd}         & Print the working directory                                 \\
	\cmd{cd dir}      & Change to \file{dir}                                        \\
	\cmd{cd}          & With no argument, go home                                   \\
	\cmd{cd -}        & Back to where you just were                                 \\
	\cmd{cd ..}       & Up one level                                                \\
	\cmd{pushd dir} / \cmd{popd} & Change directory onto a stack, and come back     \\
	\cmd{dirs -v}     & Show the directory stack                                    \\
\end{keys}

\begin{note}
	\cmd{cd} must be a shell \emph{builtin}, not a program. A child process cannot
	change its parent's working directory, so if \cmd{cd} were an ordinary command
	it would change its own directory and then exit, accomplishing nothing.
	\cmd{type cd} confirms it. The same reasoning explains \cmd{export} and
	\cmd{source}.
\end{note}

\begin{tryit}
	Run \cmd{cd /tmp}, then \cmd{cd \textasciitilde}, then \cmd{cd -} a few times.
	The last one toggles, which makes it the \key{<C-\textasciicircum{}>} of the
	shell (\autoref{sec:vim-buffercmds}) --- the fastest way to bounce between two
	directories you are working in.
\end{tryit}

\section{Listing and Inspecting}\label{sec:sh-inspect}

\begin{keys}[0.22]
	\cmd{ls}          & List the current directory                                  \\
	\cmd{ls -l}       & Long form: permissions, owner, size, date                   \\
	\cmd{ls -a}       & Include hidden files                                        \\
	\cmd{ls -lh}      & Human-readable sizes --- \texttt{4.0K} rather than \texttt{4096} \\
	\cmd{ls -lt}      & Newest first; add \cmd{-r} to reverse                       \\
	\cmd{ls -ld dir}  & The directory itself, not its contents                      \\
	\cmd{tree -L 2}   & The tree, two levels deep                                   \\
	\cmd{stat file}   & Everything the filesystem knows about one file              \\
	\cmd{file x}      & Guess what a file actually \emph{is}, ignoring its name     \\
	\cmd{du -sh dir}  & How big is this directory, in total                         \\
	\cmd{df -h}       & How full are the mounted filesystems                        \\
\end{keys}

\begin{eg}[Reading \cmd{ls -l}]
\begin{shell}
$ ls -l notes.tex
-rw-rw-r-- 1 konyi konyi 4096 Sep 14 17:23 notes.tex
\end{shell}
	Left to right: the type and permissions (\autoref{sec:sh-perms}); the link
	count; the owning user; the owning group; the size in bytes; the modification
	time; the name. The very first character is the type --- \texttt{-} for a
	regular file, \texttt{d} for a directory, \texttt{l} for a symbolic link.
\end{eg}

\begin{note}
	\cmd{file} is the honest answer to ``what is this''. Extensions are a
	convention that nothing enforces, so \cmd{file mystery.txt} reporting
	\texttt{PNG image data} is not a contradiction --- it is the point.
\end{note}

\section{Creating and Destroying}\label{sec:sh-manipulate}

\begin{keys}[0.26]
	\cmd{touch f}         & Create \file{f} if absent; update its timestamp if present \\
	\cmd{mkdir d}         & Make a directory                                           \\
	\cmd{mkdir -p a/b/c}  & Make the whole chain, and do not complain if it exists      \\
	\cmd{cp a b}          & Copy                                                       \\
	\cmd{cp -r src dst}   & Copy a directory and everything in it                      \\
	\cmd{cp -i a b}       & Ask before overwriting                                     \\
	\cmd{mv a b}          & Move --- which is also how you rename                      \\
	\cmd{rm f}            & Delete a file                                              \\
	\cmd{rm -r d}         & Delete a directory and its contents                        \\
	\cmd{rmdir d}         & Delete a directory, but only if it is empty                \\
\end{keys}

\begin{note}[There is no undo]
	\cmd{rm} does not move anything to a wastebasket. The file is gone, and the
	only recovery is your backups --- or, for a file you are tracking,
	Part~\ref{part:git}. Three habits make this survivable:
	\begin{enumerate}[(i)]
		\item Prefer \cmd{rm -i} while you are learning, or alias it so.
		\item Run \cmd{ls} on a glob before you run \cmd{rm} on it. The shell expands
		      both the same way, so what \cmd{ls} printed is exactly what \cmd{rm} will
		      delete.
		\item Never put a space in \cmd{rm -rf ./}\,\texttt{*}. \cmd{rm -rf . /}\,\texttt{*}
		      is a different and much worse command.
	\end{enumerate}
\end{note}

\begin{moral}
	\cmd{mv} is rename, copy-then-delete, and move-between-directories all at once,
	because on Unix a filename is just an entry in a directory pointing at the
	data. Renaming rewrites the entry and never touches the file.
\end{moral}

\section{Hard and Symbolic Links}\label{sec:sh-links}

That last observation has a consequence: a file's data can have more than one
name.

\begin{definition}[Hard and symbolic links]
	A \vocab{hard link} is an additional directory entry pointing at the same data.
	The file is deleted only when the last one goes. A \vocab{symbolic link} (or
	\vocab{symlink}) is a small file whose contents are a \emph{path} --- a
	signpost, which can point at something that does not exist.
\end{definition}

\begin{keys}[0.26]
	\cmd{ln target name}     & Create a hard link                                  \\
	\cmd{ln -s target name}  & Create a symbolic link --- the common case           \\
	\cmd{ls -l}              & Symlinks show as \texttt{name -> target}            \\
	\cmd{readlink -f name}   & Follow it to the real file                          \\
	\cmd{unlink name}        & Remove the link, never the target                   \\
\end{keys}

\begin{eg}
	Dotfile management is the usual reason to care. Keep the real files in a Git
	repository and symlink them into place:
\begin{shell}
$ ln -s ~/dotfiles/nvim ~/.config/nvim
$ ls -l ~/.config/nvim
lrwxrwxrwx 1 konyi konyi 22 Jul 30 ~/.config/nvim -> /home/konyi/dotfiles/nvim
\end{shell}
	Now the configuration is versioned, but Neovim still finds it exactly where it
	expects.
\end{eg}

\begin{note}
	Use an \emph{absolute} target for a symlink unless you have thought carefully.
	A relative target is resolved from the link's own directory, not from where you
	were standing when you created it, which is the cause of most broken symlinks.
\end{note}

\section{Permissions and Ownership}\label{sec:sh-perms}

Every file carries an owner, a group, and nine permission bits: read, write and
execute, for the user, the group and everyone else.

\begin{center}
	\ttfamily\large
	- \quad rw- \quad rw- \quad r-{}-
\end{center}
\begin{center}
	\small type \quad\quad user \quad\quad group \quad\quad other
\end{center}

\begin{keys}[0.18]
	\texttt{r} & Read: see the contents. On a directory, list its entries               \\
	\texttt{w} & Write: change the contents. On a directory, create or delete entries   \\
	\texttt{x} & Execute: run it. On a directory, \emph{enter} it or pass through it    \\
\end{keys}

\begin{note}
	The directory cases are the ones that surprise people. \texttt{x} without
	\texttt{r} on a directory lets you open a file inside it if you already know
	the name, but not list what is there. And deleting a file needs \texttt{w} on
	the \emph{directory}, not on the file --- which is why you can delete something
	you have no permission to read.
\end{note}

\begin{keys}[0.28]
	\cmd{chmod u+x script.sh}     & Give the owner execute permission               \\
	\cmd{chmod go-w file}         & Take write away from group and others           \\
	\cmd{chmod 644 file}          & Octal: \texttt{rw-r-{}-r-{}-}                   \\
	\cmd{chmod 755 dir}           & Octal: \texttt{rwxr-xr-x}                       \\
	\cmd{chmod -R u+rwX dir}      & Recursive; capital \texttt{X} means ``\texttt{x} only where it makes sense'' \\
	\cmd{chown user:group file}   & Change ownership --- usually needs \cmd{sudo}   \\
	\cmd{umask}                   & Which permissions are \emph{removed} from new files \\
\end{keys}

The octal form is three digits, one per class, each the sum of
$\texttt{r}=4$, $\texttt{w}=2$, $\texttt{x}=1$:

\begin{keys}[0.18]
	\texttt{600} & \texttt{rw-{}-{}-{}-{}-{}-{}-} --- private. SSH keys must be this \\
	\texttt{644} & \texttt{rw-r-{}-r-{}-} --- the default for a file                 \\
	\texttt{755} & \texttt{rwxr-xr-x} --- the default for a directory or a program   \\
	\texttt{777} & Everything for everyone. Essentially always the wrong answer      \\
\end{keys}

\begin{moral}
	When a script will not run, the question is almost never the code.
	\cmd{ls -l} it: if there is no \texttt{x}, \cmd{chmod +x} it; if there is,
	check the shebang (\autoref{sec:sh-shebang}).
\end{moral}

\section{Finding Files}\label{sec:sh-find}

\cmd{find} walks a directory tree and applies tests to everything it meets. Its
syntax is unlike every other Unix command --- long options with one dash,
arguments that read as a sentence --- and it is worth the adjustment.

\begin{center}
	\ttfamily\large
	find \{where\} \{tests\} [\{action\}]
\end{center}

\begin{keys}[0.30]
	\cmd{find . -name '*.tex'}       & By name, case-sensitively. Quote the pattern! \\
	\cmd{find . -iname '*.PDF'}      & Case-insensitively                            \\
	\cmd{find . -type f}             & Regular files only (\cmd{d} for directories)  \\
	\cmd{find . -mtime -7}           & Modified in the last seven days               \\
	\cmd{find . -size +10M}          & Larger than ten megabytes                     \\
	\cmd{find . -empty}              & Empty files and directories                   \\
	\cmd{find . -maxdepth 2}         & Do not descend further than two levels        \\
	\cmd{find . -path '*/build/*' -prune -o -print} & Skip a subtree entirely        \\
\end{keys}

Tests combine, and the default between them is \emph{and}:

\begin{shell}
$ find . -type f -name '*.log' -size +1M          # big log files
$ find . -type f -newer Makefile                  # changed since Makefile did
$ find . -name '*.aux' -delete                    # find and remove
$ find . -name '*.tex' -exec wc -l {} +           # run a command on the results
\end{shell}

\begin{note}
	In \cmd{-exec}, \texttt{\{\}} stands for each file found. Ending with
	\texttt{+} batches as many files as will fit into one invocation; ending with
	\texttt{\textbackslash;} runs the command once per file. Prefer \texttt{+}: it
	is dramatically faster on large trees, for the same reason \cmd{xargs} exists
	(\autoref{sec:sh-teexargs}).
\end{note}

\begin{moral}
	Always quote \cmd{find}'s patterns. \cmd{find . -name *.c} lets the shell
	expand the glob first (\autoref{sec:sh-glob}), and \cmd{find} then receives
	filenames where it expected a pattern --- producing either the wrong answer or
	\texttt{paths must precede expression}, which is \cmd{find}'s way of saying
	exactly this.
\end{moral}

\begin{remark}
	\cmd{fd} is a modern alternative with saner defaults: \cmd{fd '\textbackslash.tex\$'}
	does what most \cmd{find} invocations are reaching for, respects
	\file{.gitignore}, and is faster. It is not installed here, and \cmd{find} is
	on every machine you will ever log into, so learn \cmd{find} first and treat
	\cmd{fd} as a luxury.
\end{remark}
